% ---
% titre: Calcul mental
% theme: NO
% chapitre: Nombres naturels et décimaux
% sous_chapitre: PPMC et PGDC
% niveau: [9e]
% type: TH
% tags: [c-nombres-premiers, c-decomposition-facteurs-premiers, c-multiple, c-diviseur, c-ppmc,  f-commun, p-decouverte]
% difficulte: 1
% corrige: false
% ---


\header{PPMC ou PGDC ?}

\renewcommand{\arraystretch}{1.4}
\begin{tabularx}{\linewidth}{|p{3.6cm}|X|X|}
\hline
\rowcolor{themeColor!10}
\textbf{Critère} & \textbf{PPMC} & \textbf{PGDC} \\
\hline
Ce qu'on cherche & un multiple commun -- un nombre qui \emph{contient} les données & un diviseur commun -- un nombre \emph{contenu dans} les données \\
\hline
Mots-clés typiques & \og sans qu'il n'en reste \fg, \og en groupe de \fg, \og en même temps \fg, \og coïncideront \fg, \og se retrouveront \fg & \og le plus grand nombre de \fg, \og des lots identiques les plus grands possibles \fg, \og répartir équitablement \fg \\
\hline
Taille du résultat & grand (souvent $>$ que les nombres donnés) & petit (toujours $\leq$ que le plus petit des nombres donnés) \\
\hline
Résultats & on garde la puissance la plus élevée de chaque facteur premier rencontré (même s'il n'apparaît que dans un seul nombre) & on garde uniquement les facteurs premiers commun  aux deux décompositions, à la puissance la plus faible \\
\hline
\end{tabularx}
\renewcommand{\arraystretch}{1}